\section{Representability}
\label{sec:arithmetic-coding-03}

\begin{definition}[Numeralwise representation of a relation]
Let \(T\supseteq\mathsf Q\), and let \(R\subseteq\N^k\).  A formula
\(\rho_R(x_1,\ldots,x_k)\) \emph{numeralwise represents} \(R\) in \(T\) if,
for every \(\bar n\in\N^k\),
\[
\begin{aligned}
  R(\bar n)&\quad\Longrightarrow\quad
  T\derives\rho_R(\overline{\bar n}),\\
  \neg R(\bar n)&\quad\Longrightarrow\quad
  T\derives\neg\rho_R(\overline{\bar n}).
\end{aligned}
\]
The two clauses concern each standard tuple of numerals; they do not assert
that \(T\) proves a universally quantified decision statement.
\end{definition}

\begin{definition}[Representation of a function]
Let \(T\supseteq\mathsf Q\).  A formula \(\gamma_f(\bar x,y)\)
\emph{represents} a total function
\(f:\N^k\to\N\) in \(T\) if, for every \(\bar n\in\N^k\),
\[
  T\derives\forall y\,
  \bigl(\gamma_f(\overline{\bar n},y)
  \leftrightarrow y=\overline{f(\bar n)}\bigr).
\]
\end{definition}

\begin{theorem}[Representability theorem]
\label{thm:representability}
Every primitive recursive relation is numeralwise representable in
\(\mathsf Q\) by a \(\Sigma^0_1\) formula, and every primitive recursive
function is representable in \(\mathsf Q\).  Every computably enumerable
relation is definable in \(\mathcal N\) by a \(\Sigma^0_1\) formula.
\textit{Status: cited theorem.}
\end{theorem}

See \cite{Shoenfield1967} for a complete proof.

The construction represents the initial functions directly and preserves
representability through composition and primitive recursion by quantifying
over coded finite computation sequences.  The sequence-coding operations and
the check that a code satisfies the recursion equations are primitive
recursive.  For a computably enumerable relation, a finite accepting
computation supplies the existential witness.  This paragraph records the
construction used below; the cited source supplies the complete coding proof.

\begin{example}[Representation of addition]
The formula
\[
  \mathsf{Add}(x,z,y)\defeq x+z=y
\]
represents the total function \((m,n)\mapsto m+n\) in \(\mathsf Q\).  For each
standard \(m,n\), finitely many uses of the two addition axioms give
\(\mathsf Q\derives\bar m+\bar n=\overline{m+n}\).  Equality reasoning then
gives
\[
  \mathsf Q\derives\forall y\,
  \bigl(\mathsf{Add}(\bar m,\bar n,y)
  \leftrightarrow y=\overline{m+n}\bigr).
\]
This numeral-by-numeral verification uses no induction scheme.
\end{example}

\subsection*{A specified graph construction with induction in \(\PA\)}

The cited \(\mathsf Q\) theorem fixes the strength of numeralwise
representation.  For internal calculations in \(\PA\), the following
construction additionally specifies the formula and proves its defining
equations.  Its index is a finite definition of a primitive recursive
function, not an arbitrary formula defining the same function in
\(\mathcal N\).

\begin{definition}[Graph formula of a primitive recursive description]
\label{def:p05-pa-pr-graph}
A \emph{primitive recursive description} of arity \(k\) is a finite
formation tree in the initial functions, composition and primitive-recursion
clauses of \cref{sec:turing-machines-02}, with the arities recorded at each
node.  Write \(f_d\) for the numerical function denoted by such a description
\(d\).  The subscript \(d\) is a metalinguistic finite tree; it is not a
program index or an object-language variable.

Associate to \(d\) the arithmetic formula \(\Gamma_d(\bar x,y)\) as follows.
For zero, successor and projection descriptions use respectively
\(y=0\), \(y=Sx\), and \(y=x_i\).  For the composition of \(g\) with
\(h_1,\ldots,h_m\), use
\[
 \exists z_1\cdots\exists z_m\,
 \bigl(\Gamma_{h_1}(\bar x,z_1)\land\cdots\land
       \Gamma_{h_m}(\bar x,z_m)\land\Gamma_g(\bar z,y)\bigr).
\]
An empty list of intermediate variables contributes no quantifier or
conjunct; this includes composition through a nullary function.

For primitive recursion with base description \(g\) and step description
\(h\), define
\[
\begin{split}
 \Gamma_d(\bar x,n,y)\defeq\exists w\,\bigl[&
   \exists u\,(\mathsf B(w,0,u)\land\Gamma_g(\bar x,u))\\
 &\land\forall i<n\,\exists u\,\exists v\,
   (\mathsf B(w,i,u)\land\mathsf B(w,Si,v)
                  \land\Gamma_h(\bar x,i,u,v))\\
 &\land\mathsf B(w,n,y)\bigr].
\end{split}
\]
Here \(\mathsf B\) is exactly the bounded beta-entry formula of
\cref{lem:goedel-beta}.  Every bound variable introduced in these clauses is
chosen fresh from the argument variables and previously used bound variables,
in their fixed order, before substitution.  These conventions determine a
formula from the finite description.  All displayed graph names abbreviate
their recursively specified arithmetic formulas.
\end{definition}

\begin{theorem}[Primitive recursive graphs verified in \(\PA\)]
\label{thm:p05-pa-pr-graphs}
For every fixed primitive recursive description \(d\) of arity \(k\),
\(\Gamma_d\) is a \(\Sigma^0_1\) formula,
\[
 \PA\derives\forall\bar x\,\exists!y\,\Gamma_d(\bar x,y),
 \qquad
 \mathcal N\entails\Gamma_d(\overline{\bar a},\bar b)
       \quad\Longleftrightarrow\quad f_d(\bar a)=b.
\]
The uniqueness quantifier abbreviates existence together with equality of
any two outputs.  For a recursion description, \(\PA\) proves
\[
\begin{aligned}
 \Gamma_d(\bar x,0,y)&\leftrightarrow\Gamma_g(\bar x,y),\\
 \Gamma_d(\bar x,Sn,z)&\leftrightarrow
   \exists y\,(\Gamma_d(\bar x,n,y)\land\Gamma_h(\bar x,n,y,z)),
\end{aligned}
\]
with universal closure; the initial and composition equations hold by the
displayed definitions.  For each standard tuple \(\bar a\), it follows that
\[
 \PA\derives\forall y\,
   (\Gamma_d(\overline{\bar a},y)
                     \leftrightarrow y=\overline{f_d(\bar a)}).
\]
\end{theorem}
\begin{proof}
\Cref{annex:p5.internal-pr-graphs} proves the needed beta-extension and
finite-collection statements inside \(\PA\), then proves the theorem by
induction on the description.  In particular, sequence extension is proved
before it is used to justify primitive recursion.
\end{proof}

The theorem is a scheme for separately fixed finite descriptions.  It does
not assert a \(\PA\)-provably total evaluator for all primitive recursive
descriptions at once.  Nor does it identify two graph formulas from their
agreement in \(\mathcal N\).  In an application to proof coding, one must
select the description of the actual checker and prove the equations and
invariants for its proof constructors in \(\PA\).  These are further
obligations for the internal proof predicate; they are not supplied by
numeralwise representation alone.

\begin{corollary}[Lifting a recursion invariant into \(\PA\)]
\label{cor:p05-pa-recursion-invariant}
Let \(d\) be a recursion description with base \(g\) and step \(h\),
and let \(I(\bar x,n,y)\) be any arithmetic formula.  If \(\PA\)
proves the universal closures of
\[
\begin{aligned}
 \Gamma_g(\bar x,y)&\to I(\bar x,0,y),\\
 I(\bar x,n,y)\land\Gamma_h(\bar x,n,y,z)&\to I(\bar x,Sn,z),
\end{aligned}
\]
then it proves
\[
 \forall\bar x\,\forall n\,\forall y\,
   (\Gamma_d(\bar x,n,y)\to I(\bar x,n,y)).
\]
\end{corollary}
\begin{proof}
In \(\PA\), induct on \(n\) in
\(\forall y(\Gamma_d(\bar x,n,y)\to I(\bar x,n,y))\).
The zero defining equation and the first premise give the base.  For the
successor, the defining equation supplies a preceding value \(y\) and
the \(\Gamma_h\) fact.  Apply the induction hypothesis and then the
second premise.  This uses actual \(\PA\) proofs of the base and step
implications; their truth in \(\mathcal N\) alone is insufficient.
\end{proof}

\begin{example}[Specified graphs for the governing pairing and cons]
\label{ex:p05-pa-pair-cons-graphs}
The following equations are finite primitive recursive descriptions, in the
order displayed:
\[
\begin{aligned}
 A(x,0)&=x,& A(x,Sn)&=S(A(x,n)),\\
 H(0)&=0,& H(Sn)&=A(H(n),Sn),\\
 K(a,b)&=A(H(A(a,b)),b),& C(a,b)&=S(K(a,b)).
\end{aligned}
\]
The letters here name the displayed descriptions and, in the numerical
equations, their denoted functions: addition, triangular numbers,
the governing numerical pairing, and its positive-cons successor.  They
abbreviate the graphs obtained by \cref{def:p05-pa-pr-graph}, not new
primitive symbols of \(\Lang_A\).  Their graph formulas satisfy
\[
\begin{aligned}
 \PA&\derives\forall x\,\forall n\,\forall y\,
            (\Gamma_A(x,n,y)\leftrightarrow y=x+n),\\
 \PA&\derives\forall a\,\forall b\,\forall w\,
            (\Gamma_K(a,b,w)\leftrightarrow\mathsf P(a,b,w)),\\
 \PA&\derives\forall a\,\forall b\,\forall c\,
   (\Gamma_C(a,b,c)\leftrightarrow
          \exists w\,(\mathsf P(a,b,w)\land c=Sw)).
\end{aligned}
\]
For the first formula induct on \(n\): the graph's zero equation gives
\(y=x=x+0\); its step equation gives \(y=S(x+n)=x+Sn\), and these
same witnesses prove the converses.  For \(H\), induction gives
\(2H(n)=n(n+1)\) by the step identity
\(2(H(n)+n+1)=n(n+1)+2n+2=(n+1)(n+2)\).
The internal uniqueness of the triangular solution in
\cref{lem:p05-pa-pair-beta} therefore identifies \(H\) with that solution.
Composition with addition gives the displayed \(\mathsf P\) equation
for \(K\), and composition with successor gives the \(C\) equation.
Unique total graph values justify each displayed function abbreviation;
equivalently, insert their existentially quantified output variables and
graph conjuncts throughout these finite arguments.

Consequently \(\PA\) proves that a positive cons code has unique head
and tail, both strictly below that code: decode its unique predecessor by
\(\mathsf P\) and use the component bounds in
\cref{lem:p05-pa-pair-beta}.  This is an internal assertion for the specified
pairing construction, obtained from its equations rather than from the
external fact that pairing is primitive recursive.
\end{example}

\begin{theorem}[Bounded evidence for a fixed \(\Sigma^0_1\) formula in \(\PA\)]
\label{thm:p05-pa-sigma-evidence}
For every fixed \(\Sigma^0_1\) formula \(\varphi(\bar x)\), one can
construct a bounded arithmetic formula \(\mathsf{Ev}_\varphi(\bar x,e)\)
with a fresh evidence variable \(e\), such that
\[
 \PA\derives\forall\bar x\,
   (\varphi(\bar x)\leftrightarrow
                         \exists e\,\mathsf{Ev}_\varphi(\bar x,e)).
\]
The construction permits dependent existential witnesses under a marked
bounded universal.  It is an assignment to each fixed formula, not a
uniform truth predicate for arbitrary arithmetic syntax.
\end{theorem}
\begin{proof}
The final part of \cref{annex:p5.internal-pr-graphs} defines every evidence
clause and proves both directions.  The bounded-universal case uses the
internal finite-collection lemma proved there.
\end{proof}
